% =============================================================
%  Creepage distance of DC OHL insulators
%  Figures: dry-band and AC/DC arc-current sketches as inline TikZ.
%  (figures/cd_creepage_voltage.pdf and creepage_figures.py belonged to
%  the worked example, now removed from this section.)
%  Slides derived from this section:
%  presentation/design/transmission-line/creepage-distance/
% =============================================================
\section{Creepage Distance}
\label{sec:creepage}

\subsection{Definition and reference voltage}
\label{sec:creepage-definition}

The creepage distance is the shortest distance along the surface of an
insulator between the two conductive parts that normally carry the operating
voltage. On a DC overhead line it is set by pollution, not by overvoltage: a
wetted pollution layer conducts, and an insulator with too little creepage
can flash over at normal operating voltage.

The creepage is specified as a unified specific creepage distance (USCD) in
mm/kV, referred to the pole-to-earth DC operating voltage~[1]. The total
creepage of an insulator string follows as
\begin{equation}
  L = \mathrm{USCD} \cdot U_{dpe} \cdot C_{a}
  \label{eq:creepage}
\end{equation}
with $L$ the total creepage in mm, $U_{dpe}$ the pole-to-earth DC operating
voltage in kV and $C_{a}$ the altitude correction of
Section~\ref{sec:creepage-altitude}. There is no
single recommended USCD: it follows from the site pollution severity and the
insulator type~[1], [2].

\subsection{Pollution flashover and the dry-band arc}
\label{sec:creepage-dryband}

A pollution flashover is caused by wet pollution. CIGRE Working Group B4.45
describes the process in detail~[3]. Two forms of pollution are
distinguished~[4]: inland, the pollutants settle as dry dust and are wetted
later by rain, fog or condensation (Type A); near the coast or at sea they
arrive already wet, for example as salt fog (Type B). On an insulator with a
hydrophilic surface the flashover then develops in six stages~[3]:

\begin{enumerate}[itemsep=0.2em]
  \item \textbf{Pollution accumulation.} Pollutants from the air settle on
        the surface. How much accumulates depends on the amount in the air,
        on the electric field through electrostatic forces, and on the shed
        profile. The deposit may contain several kinds of salt.
  \item \textbf{Wetting and dissolution.} Rain, fog, condensation or
        adsorption from humid air brings water to the deposit and dissolves
        at least part of it, so the layer becomes conductive. Different
        pollutants dissolve, and absorb water, at different rates.
  \item \textbf{Leakage current.} The wet layer conducts, and a leakage
        current of a few mA to a few hundred mA flows along the surface.
  \item \textbf{Dry-band formation.} The current heats the layer. Neither the
        pollution nor the water is evenly distributed, and the current density
        is highest near the pin and cap of a disc or on the narrow core of a
        long-rod or composite insulator, so some paths dry faster than others
        and non-conductive dry bands form.
  \item \textbf{Partial arcs.} Nearly the full voltage of the section now
        stands across dry bands a few mm to a few cm wide. The field exceeds
        the strength of air, and partial arcs bridge the bands, burning in
        series with the remaining wet layer.
  \item \textbf{Total flashover.} In the dynamic interplay of the dry-band
        arcs, one arc may extend across the whole insulator.
\end{enumerate}

Two points follow~[3]. A flashover needs both pollution and wetting: without
pollution, or without the pollution becoming wet, none occurs. And both the
accumulation and the flashover itself follow different patterns under AC and
DC, Section~\ref{sec:creepage-dc}.

The arc and the wet layer share the applied voltage. For an arc of length $x$
on a creepage path of length $L$, the model of Obenaus~[5] gives
\begin{equation}
  U = x \cdot A \cdot I^{-n} + R' \cdot (L - x) \cdot I
  \label{eq:obenaus}
\end{equation}
where the first term is the arc voltage, whose gradient falls with rising
current $I$ ($A$ and $n$ are arc constants), and the second term is the
voltage across the remaining wet layer of resistance $R'$ per unit length.
Figure~\ref{fig:dryband}(a) shows the arrangement.

If the arc needs a lower voltage gradient than the wet layer it replaces, it
elongates along the surface~[6]. The current then rises, the arc gradient
falls further, and the arc runs away until it bridges the whole insulator ---
the flashover. Otherwise it stalls, wanders and extinguishes, and appears
only as surface sparking and leakage-current pulses. The creepage distance
has to be long enough that the runaway condition is never reached at
operating voltage.

\subsection{Why DC needs more creepage than AC}
\label{sec:creepage-dc}

The same site needs a markedly higher USCD on DC than on AC, for three
reasons:

\begin{itemize}[itemsep=0.2em]
  \item \textbf{More pollution.} The steady DC field attracts charged
        particles, so a DC insulator collects roughly 1.5--2 times the
        pollution of an AC insulator at the same site~[2].
  \item \textbf{No current zero.} An AC dry-band arc extinguishes at every
        current zero, 100 or 120 times per second, and has to re-ignite across
        a band that cools in the meantime. A DC arc burns continuously and can
        keep elongating without interruption, Figure~\ref{fig:dryband}(b).
  \item \textbf{Ionic migration.} Under DC, ions migrate through glass and
        porcelain and corrode the pin; DC insulators need DC-type glass and a
        zinc sleeve on the pin.
\end{itemize}

\begin{figure}[!ht]
  \centering
  \begin{tikzpicture}[font=\small, >=Latex]
    % ---- (a) dry band and arc on the unrolled creepage path ----
    \node[font=\small\bfseries] at (-1.2,1.6) {(a)};
    % electrodes
    \fill[atBorderStrong] (-0.4,-0.35) rectangle (0,0.55);
    \fill[atBorderStrong] (11,-0.35) rectangle (11.4,0.55);
    \node[above, align=center, font=\footnotesize] at (-0.2,0.6) {pole\\$U_{dpe}$};
    \node[above, align=center, font=\footnotesize] at (11.2,0.6) {earth};
    % insulator surface
    \fill[atBgMuted] (0,-0.35) rectangle (11,0);
    % wet layer left and right of the dry band
    \fill[atInfo!45] (0,0) rectangle (2.4,0.18);
    \fill[atInfo!45] (3.6,0) rectangle (11,0.18);
    % dry band
    \draw[atFgMuted, densely dashed] (2.4,0) -- (2.4,0.18);
    \draw[atFgMuted, densely dashed] (3.6,0) -- (3.6,0.18);
    \node[below, font=\footnotesize] at (3.0,-0.35) {dry band};
    % arc
    \draw[atDanger, very thick, decorate,
          decoration={snake, amplitude=0.6mm, segment length=2.5mm}]
         (2.3,0.18) .. controls (2.6,1.1) and (3.4,1.1) .. (3.7,0.18);
    \node[atDanger, font=\footnotesize] at (3.0,1.25) {arc};
    % leakage current
    \draw[->, atFg] (5.0,0.42) -- (8.0,0.42) node[midway, above, font=\footnotesize] {leakage current $I$};
    \node[font=\footnotesize, atFgMuted] at (7.3,-0.18) {wet pollution layer, $R'$ per unit length};
    % dimensions
    \draw[|<->|, atFgMuted] (2.4,-0.85) -- (3.6,-0.85) node[midway, below, font=\footnotesize] {$x$};
    \draw[|<->|, atFgMuted] (3.6,-0.85) -- (11,-0.85) node[midway, below, font=\footnotesize] {$L - x$};

    % ---- (b) arc current, AC versus DC ----
    \begin{scope}[yshift=-3.4cm]
      \node[font=\small\bfseries] at (-1.2,1.3) {(b)};
      % AC
      \draw[->, atFgMuted] (0,0) -- (5.2,0) node[right, font=\footnotesize] {$t$};
      \draw[->, atFgMuted] (0,-1.0) -- (0,1.2) node[above, font=\footnotesize] {$i_{arc}$};
      \draw[atInfo, thick, domain=0:4.8, samples=120] plot (\x, {0.85*sin(\x*150)});
      \foreach \z in {1.2,2.4,3.6,4.8} \fill[atDanger] (\z,0) circle (1.6pt);
      \node[align=center, font=\footnotesize] at (2.5,-1.45)
           {AC: arc extinguishes at every current zero\\and must re-ignite};
      % DC
      \begin{scope}[xshift=6.2cm]
        \draw[->, atFgMuted] (0,0) -- (4.8,0) node[right, font=\footnotesize] {$t$};
        \draw[->, atFgMuted] (0,-1.0) -- (0,1.2) node[above, font=\footnotesize] {$i_{arc}$};
        \draw[atAccentFg, thick] (0,0.7) -- (4.6,0.7);
        \node[align=center, font=\footnotesize] at (2.4,-1.45)
             {DC: no current zero,\\the arc burns and elongates};
      \end{scope}
    \end{scope}
  \end{tikzpicture}
  \caption{(a) Dry-band arc on the unrolled creepage path: a partial arc of
  length $x$ bridges the dry band and burns in series with the wet layer of
  length $L - x$, Equation~\eqref{eq:obenaus}. (b) Arc current on AC and DC.}
  \label{fig:dryband}
\end{figure}

Table~\ref{tab:acdc} summarises the differences. The composite (silicone)
insulator counters the mechanism at its root: its water-repellent surface
breaks the wet film into separate droplets, keeps the leakage current low
and prevents dry bands from forming. Alternating large and small sheds with
sufficient spacing also stop arcs from bridging between sheds and shortening
the creepage path.

\begin{table}[!ht]
  \centering
  \begin{tabularx}{\linewidth}{@{}>{\raggedright\arraybackslash}p{0.27\linewidth}
                                 >{\raggedright\arraybackslash}X
                                 >{\raggedright\arraybackslash}X@{}}
    \toprule
    \multicolumn{1}{c}{\textbf{Aspect}} & \textbf{AC} & \textbf{DC} \\
    \midrule
    Pollution build-up      & wind and gravity      & wind and gravity plus electrostatic attraction \\
    Dry-band arc            & extinguishes at every current zero & no current zero, continuous \\
    Ionic migration in glass and porcelain & negligible & DC-type glass, zinc sleeve on the pin \\
    Relative USCD           & 1                     & about 1.5--2 \\
    \bottomrule
  \end{tabularx}
  \caption{Pollution performance of AC and DC insulators.}
  \label{tab:acdc}
\end{table}

\subsection{Site pollution severity: ESDD and NSDD}
\label{sec:creepage-esdd}

The pollution on an insulator has two parts. The \emph{soluble deposit} is
the salt --- sea salt, industrial sulfates and nitrates, fertiliser, road
de-icing salt, or the sodium, calcium and magnesium salts in desert and
dry-lake dust --- that dissolves into ions when wetted and makes the layer
conduct. The \emph{non-soluble deposit} --- clay, sand, soot --- does not
conduct but holds water and keeps the soluble film wet for longer.

\begin{description}[itemsep=0.2em]
  \item[ESDD] Equivalent salt deposit density: the mass of NaCl per surface
        area that, dissolved in water, gives the same conductivity as the
        soluble deposit washed off the insulator, in mg/cm$^2$~[4]. Converting
        to NaCl puts every salt on one scale, since different salts conduct
        differently per gram.
  \item[NSDD] Non-soluble deposit density: the mass of the non-soluble residue
        per surface area, in mg/cm$^2$~[4].
\end{description}

Both are measured on reference insulators exposed on site for a year or more.
A defined area $A$ in cm$^2$ is washed with a volume $V$ in cm$^3$ of
demineralised water. The conductivity $\sigma_{\theta}$ of the wash water at
temperature $\theta$ is corrected to 20\,$^{\circ}$C with the temperature
factor $b$ of~[7], and converted into the salinity $S_{a}$ in kg/m$^3$
(equal to mg/cm$^3$):
\begin{align}
  \sigma_{20} &= \sigma_{\theta} \cdot \left[1 - b \cdot (\theta - 20\,^{\circ}\mathrm{C})\right] \\
  S_{a} &= \left(5.7 \cdot \sigma_{20}\right)^{1.03} \\
  \mathrm{ESDD} &= \frac{S_{a} \cdot V}{A}
\end{align}
with $\sigma$ in S/m. The wash water is then filtered; with $W_{i}$ and
$W_{f}$ the dry mass of the filter in g before and after,
\begin{equation}
  \mathrm{NSDD} = \frac{1000 \cdot (W_{f} - W_{i})}{A}
\end{equation}

Slowly soluble salts such as gypsum (CaSO$_4$), common in cement and some
desert dusts, dissolve only partly; the measured ESDD then depends on the
water volume and stirring time, so the wash procedure has to be fixed and
stated with the results.

Together, ESDD and NSDD give the SPS class of the site, from very light to
very heavy~[4]. For DC, the SPS must be measured on DC-energised insulators,
or estimated from AC or non-energised measurements with a DC/AC accumulation
ratio of about 1.5--2~[1], [2].

The flashover strength per unit creepage falls with both quantities as a
power law,
\begin{equation}
  E_{f} = \frac{U_{f}}{L} = A_{f} \cdot \mathrm{ESDD}^{-a} \cdot \mathrm{NSDD}^{-b_{n}}
\end{equation}
with $A_{f}$ a constant of the insulator type. The required USCD therefore
scales as
\begin{equation}
  \mathrm{USCD}_{2} = \mathrm{USCD}_{1}
    \cdot \left(\frac{\mathrm{ESDD}_{2}}{\mathrm{ESDD}_{1}}\right)^{a}
    \cdot \left(\frac{\mathrm{NSDD}_{2}}{\mathrm{NSDD}_{1}}\right)^{b_{n}}
\end{equation}
Typical exponents are $a \approx 0.2$--0.4 and $b_{n} \approx 0.1$--0.2, so
the salt dominates: doubling the ESDD with $a = 0.3$ needs
$2^{0.3} \approx 1.23$, i.e.\ 23\,\% more creepage. The exponents depend on
insulator type and profile and are taken from the supplier's DC pollution
tests or published data~[2].

\subsection{Indicative USCD values}
\label{sec:creepage-uscd}

Table~\ref{tab:uscd} gives indicative USCD ranges for DC OHL insulators.
Composite insulators need roughly 20--30\,\% less creepage than ceramic or
glass, which is why most recent HVDC lines use composite long-rod
insulators. The ranges are a starting point; the design value comes from the
measured SPS.

\begin{table}[!ht]
  \centering
  \begin{tabular}{@{}l r r@{}}
    \toprule
    \multicolumn{1}{c}{\textbf{Site pollution}}
      & \multicolumn{1}{c}{\textbf{Ceramic / glass}}
      & \multicolumn{1}{c}{\textbf{Composite}} \\
    \midrule
    Light      & 30--40\,mm/kV   & 25--30\,mm/kV \\
    Medium     & 40--50\,mm/kV   & 30--38\,mm/kV \\
    Heavy      & 50--60\,mm/kV   & 38--45\,mm/kV \\
    Very heavy & mostly too long & 45--55\,mm/kV \\
    \bottomrule
  \end{tabular}
  \caption{Indicative USCD for DC OHL insulators~[1], [2].}
  \label{tab:uscd}
\end{table}

\subsection{Altitude correction}
\label{sec:creepage-altitude}

At lower air pressure the dry-band arc needs a lower gradient, so the
pollution flashover voltage falls and more creepage is needed. Pollution
tests show a power law in the pressure ratio, and the pressure falls
exponentially with the altitude $H$~[8], [9]:
\begin{equation}
  U_{f} = U_{f0} \cdot \left(\frac{p}{p_{0}}\right)^{m}
  \qquad
  \frac{p}{p_{0}} = e^{-H/8150\,\mathrm{m}}
\end{equation}
with $U_{f}$ and $U_{f0}$ the pollution flashover voltage at pressure $p$ and
at sea level $p_{0}$. Since the required creepage is inversely proportional to
the flashover voltage, the creepage has to be multiplied by
\begin{equation}
  C_{a} = e^{\,m \cdot \frac{H}{8150\,\mathrm{m}}}
  \label{eq:altitude-full}
\end{equation}
This is the correction factor of IEC TS 60815-4~[1] and TB 518~[2],
recommended for DC polluted insulators of both ceramic and
polymeric material. CIGRE Working Group D1.44 reviews it in TB 705~[8],
where the exponent is called $n$; it is written $m$ here to keep it apart
from the arc constant $n$ of Equation~\eqref{eq:obenaus}.

\subsubsection{Exponent $m$.}
$m$ is empirical: the slope of $\ln U_{f}$ over $\ln (p/p_{0})$ from
pollution tests at several pressures in a climate chamber or at high-altitude
laboratories. It depends on the insulator type and material, the shed
profile, the pollution level and the voltage type and polarity, and lies
between 0 and 1 because only the arc depends on pressure, not the resistance
of the wet layer. For DC, IEC TS 60815-4 and TB 518 suggest $m = 0.35$ until
further information is available, a value still under consideration by
CIGRE and IEC~[8]. Since the correction may depend on the insulator
characteristics and the pollution severity, it is applied per candidate
insulator~[1]. For comparison, IEC 60071-2 gives $m = 0.5$ for standard and
up to 0.8 for anti-fog cap-and-pin insulators under AC~[9].


\paragraph{Reference height.}
Two ways of counting the height are in use~[8]:

\begin{enumerate}[label=(\alph*), itemsep=0.2em]
  \item \textbf{Full height above sea level}, as in IEC 60071-2 and
        Equation~\eqref{eq:altitude-full}. The USCD values are taken as
        sea-level values, and every site above sea level is corrected.
  \item \textbf{Height above 1000\,m}, as in some industry practice. The USCD
        values are taken as valid up to 1000\,m, and $H$ is replaced by the
        height above 1000\,m:
        \begin{equation}
          C_{a} =
          \begin{cases}
            1 & H \le 1000\,\mathrm{m}\\[0.3em]
            e^{\,m \cdot \frac{H - 1000\,\mathrm{m}}{8150\,\mathrm{m}}} & H > 1000\,\mathrm{m}
          \end{cases}
          \label{eq:altitude-1000}
        \end{equation}
\end{enumerate}

Variant (a) is the more conservative: it already adds 4.4\,\% at 1000\,m, and
above 1000\,m it stays higher than variant (b) by the constant factor
$e^{0.35 \cdot 1000/8150} = 1.044$. The specification has to state which
variant applies, together with $m$. Table~\ref{tab:altitude} compares both
for $m = 0.35$, with the resulting creepage of a 400\,kV ceramic string at
medium pollution ($\mathrm{USCD} = 45\,\mathrm{mm/kV}$, so
$L = 18.0\,\mathrm{m}$ before correction).

\begin{table}[!ht]
  \centering
  \begin{tabular}{@{}r r r r r@{}}
    \toprule
    & \multicolumn{2}{c}{\textbf{(a) Full height}}
    & \multicolumn{2}{c}{\textbf{(b) Above 1000\,m}} \\
    \cmidrule(lr){2-3}\cmidrule(l){4-5}
    \multicolumn{1}{c}{\textbf{Altitude $H$}}
      & \multicolumn{1}{c}{\textbf{$C_{a}$}} & \multicolumn{1}{c}{\textbf{$L_{cer}$}}
      & \multicolumn{1}{c}{\textbf{$C_{a}$}} & \multicolumn{1}{c}{\textbf{$L_{cer}$}} \\
    \midrule
    1000\,m & 1.044 & 18.8\,m & 1.000 & 18.0\,m \\
    1500\,m & 1.067 & 19.2\,m & 1.022 & 18.4\,m \\
    2000\,m & 1.090 & 19.6\,m & 1.044 & 18.8\,m \\
    2500\,m & 1.113 & 20.0\,m & 1.067 & 19.2\,m \\
    3000\,m & 1.138 & 20.5\,m & 1.090 & 19.6\,m \\
    \bottomrule
  \end{tabular}
  \caption{Altitude correction $C_{a}$ for $m = 0.35$ and the resulting
  creepage of the 400\,kV ceramic string ($L = 18.0\,\mathrm{m}$ uncorrected).}
  \label{tab:altitude}
\end{table}

\FloatBarrier
TB 417 treats the altitude correction of the external insulation of outdoor
DC yard equipment in the converter station~[3, Section~7.1.8.3]:

\begin{itemize}[itemsep=0.2em]
  \item The air clearances of the DC switchyard are corrected for altitudes
        above 1000\,m. The air clearances of the DC yard apparatus and the
        applicable test voltages are corrected above 1000\,m according to the
        related equipment standards.
  \item For the creepage distance there is no generally accepted rule on
        altitude correction in the current international standards.
  \item Although the lower air density at high altitude may affect the
        insulation capability along the equipment surface, it is usual
        industrial practice not to correct creepage distances for altitude,
        in AC and DC substations alike.
  \item TB 417 considers it acceptable to apply the same design method to
        800\,kV DC equipment, i.e.\ no altitude correction of the creepage
        distance, as long as the specific creepage distance is selected
        conservatively and composite material is used for the equipment
        housing.
\end{itemize}

These statements concern station equipment. For the insulators of the DC
overhead line, IEC TS 60815-4 and TB 518 recommend the correction
$C_{a}$~[1], [2], and this report applies it with
Equation~\eqref{eq:altitude-full} or~\eqref{eq:altitude-1000}.


\subsection{Design procedure}
\label{sec:creepage-procedure}

\begin{enumerate}[itemsep=0.2em]
  \item Determine the SPS along the route: ESDD and NSDD from test insulators
        exposed under DC voltage, or from service experience~[1], [2].
  \item Convert the SPS into a USCD with the IEC TS 60815-4 and CIGRE curves,
        or with the supplier's DC pollution test data for the chosen
        insulator profile.
  \item Apply the altitude correction $C_{a}$,
        Equation~\eqref{eq:altitude-full} or \eqref{eq:altitude-1000}, with
        $m$ and the height convention stated in the specification, and the
        corrections for insulator diameter and profile.
  \item Check the resulting string against the switching and lightning
        impulse requirements; the longer of the two governs.
\end{enumerate}
